\chapter{Literature Review}
\label{ch:litreview}

\section{Theoretical Background}
\label{sec:theory-background}

\subsection{The Standard Model Extension}
\label{sec:sme}

The Standard Model Extension (SME) is constructed as the most general
effective field theory Lagrangian obtainable by adding to the Standard
Model (and General Relativity) every observer-scalar term formed by
contracting Standard Model field operators with coefficients carrying free
Lorentz indices (Colladay and Kostelecký, 1997, 1998). A distinction
central to the framework is between two kinds of Lorentz transformation.
\emph{Observer} transformations rotate or boost the coordinate frame used
to describe a system without touching the system itself; the SME is
constructed to respect observer Lorentz invariance exactly. \emph{Particle}
transformations rotate or boost a localised particle or field
configuration while holding the background coefficients fixed, and it is
only under this class of transformation that the SME's background
coefficients act as fixed, non-dynamical tensors that break the symmetry.
This distinction is what allows the SME to remain a coordinate-independent,
internally consistent field theory while still describing genuine,
in-principle observable violations of Lorentz invariance.

The SME Lagrangian is written schematically as
\begin{equation}
  \mathcal{L}_{\mathrm{SME}} = \mathcal{L}_{\mathrm{SM}} + \delta\mathcal{L}_{\mathrm{LV}},
\end{equation}
where $\mathcal L_{\mathrm{SM}}$ is the ordinary Standard Model Lagrangian
and $\delta\mathcal L_{\mathrm{LV}}$ collects every observer-Lorentz-scalar,
particle-Lorentz-violating operator that can be built from Standard Model
fields. This thesis works within the \emph{minimal} SME, the subset of
$\delta\mathcal L_{\mathrm{LV}}$ of renormalisable mass dimension
($\leq4$) in the fermion sector, where existing spin-dependent
exotic-interaction experiments place their bounds.

For a single Dirac fermion $\psi$ of mass $m$, the general minimal SME
fermion-sector Lagrangian is
\begin{equation}
  \mathcal L_{\mathrm{fermion}}
  = \frac{i}{2}\bar\psi\,\Gamma^\nu\overleftrightarrow{D}_\nu\,\psi
  - \bar\psi\,M\,\psi, \label{eq:sme-general}
\end{equation}
where $D_\nu$ is the gauge-covariant derivative,
$\overleftrightarrow{D}_\nu \equiv D_\nu - \overleftarrow{D}_\nu$, and the
Lorentz-violating physics is packaged entirely into two composite
matrix-valued objects,
\begin{align}
  \Gamma^\nu &= \gamma^\nu + c^{\mu\nu}\gamma_\mu + d^{\mu\nu}\gamma_5\gamma_\mu
               + e^\nu + if^\nu\gamma_5 + \tfrac12 g^{\lambda\mu\nu}\sigma_{\lambda\mu},
               \label{eq:Gamma-general} \\
  M &= m + a_\mu\gamma^\mu + b_\mu\gamma_5\gamma^\mu + \tfrac12 H_{\mu\nu}\sigma^{\mu\nu}.
      \label{eq:M-general}
\end{align}
$\Gamma^\nu$ modifies the kinetic (derivative-coupling) sector, while $M$
modifies the mass sector. Every coefficient in Eqs.~\eqref{eq:Gamma-general}
and \eqref{eq:M-general} is a fixed background tensor, not a dynamical
field, and each is in principle independent for every Standard Model fermion
species. Setting all coefficients to zero except $m$ recovers the ordinary
Dirac Lagrangian. For the leading-order (dimension $\leq4$) subset used
throughout this thesis, Eq.~\eqref{eq:sme-general} reduces to the compact
form
\begin{multline}
  \mathcal L_{\mathrm{fermion}}
  = \bar\psi\left(i\gamma^\mu D_\mu - m\right)\psi
  - a_\mu\,\bar\psi\gamma^\mu\psi
  - b_\mu\,\bar\psi\gamma_5\gamma^\mu\psi \\
  - \tfrac12 H_{\mu\nu}\,\bar\psi\sigma^{\mu\nu}\psi
  - c_{\mu\nu}\,\bar\psi\gamma^\mu D^\nu\psi
  - d_{\mu\nu}\,\bar\psi\gamma_5\gamma^\mu D^\nu\psi,
  \label{eq:sme-minimal}
\end{multline}
where $D_\mu$ is the gauge-covariant derivative and
$\sigma^{\mu\nu}=\tfrac{i}{2}[\gamma^\mu,\gamma^\nu]$. Note the operator
ordering in the $b_\mu$ term, $\gamma_5\gamma^\mu$ rather than
$\gamma^\mu\gamma_5$: since $\{\gamma_5,\gamma^\mu\}=0$, the two orderings
differ by an overall sign, so this convention fixes the sign of every
non-relativistic result derived from $b_\mu$ later in this thesis. Table
\ref{tab:sme-classification} classifies every coefficient in
Eq.~\eqref{eq:sme-minimal}, together with the dimension-5 coefficients
$e_\mu,f_\mu$, by CPT parity and leading spin-dependent role.

\begin{table}[h]
\centering
\begin{tabular}{@{}llc>{\raggedright\arraybackslash}p{6.5cm}@{}}
\toprule
Coefficient & CPT & Mass dim.\ & Leading spin-dependent role \\
\midrule
$a_\mu$        & Odd  & 4 & None at leading order; velocity-dependent at $O(1/m)$ \\
$b_\mu$        & Odd  & 4 & Zeeman-like spin coupling, $V_2$ \\
$c_{\mu\nu}$   & Even & 4 & Direction-dependent kinetic modification; subleading $V_2$-type \\
$d_{\mu\nu}$   & Even & 4 & Spin-velocity and mass-enhanced couplings, $V_2,V_7,V_8$ \\
$e_\mu$        & Odd  & 5 & Velocity-dependent, contributes to $V_8$ and higher \\
$f_\mu$        & Odd  & 5 & Velocity-dependent, contributes to $V_8$ and higher \\
$g_{\lambda\mu\nu}$ & Odd & 5 & Not treated in this thesis \\
$H_{\mu\nu}$   & Even & 4 & Anisotropic spin coupling, $V_3,V_7$ \\
\bottomrule
\end{tabular}
\caption{Classification of minimal (and near-minimal) SME fermion-sector
coefficients (Colladay and Kostelecký, 1998; Kostelecký and Lane, 1999).}
\label{tab:sme-classification}
\end{table}

The CPT parity of a coefficient follows from how its associated fermion
bilinear $\bar\psi\Gamma\psi$ transforms under the combined charge
conjugation, parity, and time-reversal operation,
\begin{equation}
  \mathrm{CPT}\!\left[\bar\psi\,\Gamma\,\psi\right]
  = \eta_\Gamma\,\bar\psi\,\Gamma\,\psi,
  \qquad
  \eta_\Gamma = \begin{cases}
    +1 & \Gamma \in \{\mathbb 1,\; \gamma^\mu,\; \sigma^{\mu\nu}\} \\
    -1 & \Gamma \in \{\gamma_5,\; \gamma_5\gamma^\mu\}
  \end{cases},
\end{equation}
and the corresponding non-relativistic Hamiltonians obey the general rule
\begin{equation}
  H_{NR}^{\bar f}\!\left(X^{\text{CPT-odd}}\right)  = -H_{NR}^{f}\!\left(X^{\text{CPT-odd}}\right),
  \qquad
  H_{NR}^{\bar f}\!\left(X^{\text{CPT-even}}\right) = +H_{NR}^{f}\!\left(X^{\text{CPT-even}}\right).
  \label{eq:cpt-rule-general}
\end{equation}

\subsection{The Dobrescu--Mocioiu Classification of Exotic Potentials}
\label{sec:dm}

Experimental searches for new long-range forces are conducted in the
non-relativistic laboratory frame, not in the relativistic language of the
SME. Dobrescu and Mocioiu (2006) closed this gap by enumerating, from first
principles, every non-relativistic potential that can arise from the
exchange of a single boson of arbitrary spin and parity between two
spin-$\tfrac12$ fermions. The result is a complete basis of sixteen
potentials,
\begin{equation}
  V(\mathbf r) = \sum_{i=1}^{16} V_i(\mathbf r),
\end{equation}
parametrised by coupling constants $g_s,g_p,g_V,g_A$ (scalar, pseudoscalar,
vector, axial-vector). With $y(r)=e^{-r/\lambda}/(4\pi)$,
$\lambda=1/m_\phi$ the mediator's Compton wavelength, $\hat{\mathbf r}$
pointing from particle 2 to particle 1, and $\mathbf v$ the relative
velocity of the pair, the sixteen potentials are (Cong \textit{et al.},
2025; Dobrescu and Mocioiu, 2006):

{\small
\begin{align}
  V_1 &= \frac1r\,y(r) \label{eq:v1}\\
  V_2 &= \frac1r\,(\boldsymbol\sigma_1\cdot\boldsymbol\sigma_2)\,y(r) \label{eq:v2}\\
  V_3 &= \frac{e^{-r/\lambda}}{4\pi m^2}\left\{\boldsymbol\sigma_1\cdot\boldsymbol\sigma_2
    \left[\frac1{r^3}+\frac1{\lambda r^2}\right]
    -(\boldsymbol\sigma_1\cdot\hat{\mathbf r})(\boldsymbol\sigma_2\cdot\hat{\mathbf r})
    \left[\frac3{r^3}+\frac3{\lambda r^2}+\frac1{\lambda^2r}\right]\right\} \label{eq:v3}\\
  V_{4,5} &= -\frac{1}{2mr^2}(\boldsymbol\sigma_1\pm\boldsymbol\sigma_2)\cdot(\mathbf v\times\hat{\mathbf r})
    \left(1-r\frac{d}{dr}\right)y(r) \label{eq:v45}\\
  V_{6,7} &= -\frac{1}{2mr^2}\Big[(\boldsymbol\sigma_1\cdot\mathbf v)(\boldsymbol\sigma_2\cdot\hat{\mathbf r})
    \pm(\boldsymbol\sigma_1\cdot\hat{\mathbf r})(\boldsymbol\sigma_2\cdot\mathbf v)\Big]
    \left(1-r\frac{d}{dr}\right)y(r) \label{eq:v67}\\
  V_8 &= \frac1r(\boldsymbol\sigma_1\cdot\mathbf v)(\boldsymbol\sigma_2\cdot\mathbf v)\,y(r) \label{eq:v8}\\
  V_{9,10} &= -\frac{1}{2mr^2}(\boldsymbol\sigma_1\pm\boldsymbol\sigma_2)\cdot\hat{\mathbf r}
    \left(1-r\frac{d}{dr}\right)y(r) \label{eq:v910}\\
  V_{11} &= -\frac{1}{mr^2}(\boldsymbol\sigma_1\times\boldsymbol\sigma_2)\cdot\hat{\mathbf r}
    \left(1-r\frac{d}{dr}\right)y(r) \label{eq:v11}\\
  V_{12,13} &= \frac{1}{2r}(\boldsymbol\sigma_1\pm\boldsymbol\sigma_2)\cdot\mathbf v\,y(r) \label{eq:v1213}\\
  V_{14} &= \frac1r(\boldsymbol\sigma_1\times\boldsymbol\sigma_2)\cdot\mathbf v\,y(r) \label{eq:v14}\\
  V_{15} &= -\frac{3}{2m^2r^3}\Big\{[\boldsymbol\sigma_1\cdot(\mathbf v\times\hat{\mathbf r})](\boldsymbol\sigma_2\cdot\hat{\mathbf r})
    +(\boldsymbol\sigma_1\cdot\hat{\mathbf r})[\boldsymbol\sigma_2\cdot(\mathbf v\times\hat{\mathbf r})]\Big\}
    \left(1-r\frac{d}{dr}+\frac{r^2}{3}\frac{d^2}{dr^2}\right)y(r) \label{eq:v15}\\
  V_{16} &= -\frac{1}{2mr^2}\Big\{[\boldsymbol\sigma_1\cdot(\mathbf v\times\hat{\mathbf r})](\boldsymbol\sigma_2\cdot\mathbf v)
    +(\boldsymbol\sigma_1\cdot\mathbf v)[\boldsymbol\sigma_2\cdot(\mathbf v\times\hat{\mathbf r})]\Big\}
    \left(1-r\frac{d}{dr}\right)y(r) \label{eq:v16}
\end{align}
}

Table~\ref{tab:dm-classification} groups the sixteen potentials by
structural class, the organising principle Dobrescu and Mocioiu use to
enumerate them exhaustively, and the property that determines which SME
coefficient(s) can generate a given $V_i$ (Chapter~4).

\begin{table}[h]
\centering
\begin{tabular}{@{}ll@{}}
\toprule
Structural class & Potentials \\
\midrule
Monopole--monopole (spin-independent) & $V_1$ \\
Monopole--dipole & $V_9, V_{10}$ \\
Dipole--dipole (spin--spin) & $V_2, V_3$ \\
Spin--velocity & $V_4, V_5, V_6, V_7$ \\
Velocity--velocity & $V_8$ \\
Spin cross-product & $V_{11}$ \\
Single-spin--velocity & $V_{12}, V_{13}, V_{14}$ \\
Higher-order spin--velocity & $V_{15}, V_{16}$ \\
\bottomrule
\end{tabular}
\caption{Structural classification of the Dobrescu--Mocioiu potentials.}
\label{tab:dm-classification}
\end{table}

\subsection{The Foldy--Wouthuysen Transformation}
\label{sec:fw}

The Foldy--Wouthuysen transformation is selected over alternative
non-relativistic reduction methods for two reasons specific to this thesis.
The Klein--Gordon equation is excluded outright: it describes spin-0 fields
and carries no spinor structure, making it incapable of generating the
spin-operator combinations $\bm\sigma_1\cdot\bm\sigma_2$,
$\bm\sigma\cdot\hat r$, and $\bm\sigma\cdot\mathbf v$ that define the
$V_1$--$V_{16}$ basis. Working directly with the full Dirac equation, while
formally correct, leaves the result in a relativistic form incompatible
with the non-relativistic experimental observables the DM potentials are
written for; the FW transformation resolves this by producing a
block-diagonal Hamiltonian in which the upper block is immediately readable
as a non-relativistic Pauli-operator expression, order by order in $1/m$,
ready for direct matching against $V_i$.

The Foldy--Wouthuysen (FW) transformation (Foldy and Wouthuysen, 1950) is
the standard technique for systematically extracting the non-relativistic
limit of a relativistic Dirac Hamiltonian, order by order in $1/m$, while
keeping every step exact and unitary. It decouples the upper (particle) and
lower (antiparticle) blocks of the Dirac spinor by eliminating, order by
order, the operators that mix the two blocks. Writing
$H = \beta m + \varepsilon_{\text{even}} + \varepsilon_{\text{odd}}$, where
$\varepsilon_{\text{even}}$ commutes with $\beta=\gamma^0$ and
$\varepsilon_{\text{odd}}$ anticommutes with it, the transformed Hamiltonian
expands to order $1/m$ as
\begin{equation}
  H' \approx \beta m + \varepsilon_{\text{even}}
  + \frac{\beta\varepsilon_{\text{odd}}^2}{2m}
  - \frac{[\varepsilon_{\text{odd}},\varepsilon_{\text{even}}]}{4m}
  + O(m^{-2}). \label{eq:FWexpansion}
\end{equation}
An even piece of the original perturbation survives into the
non-relativistic Hamiltonian unchanged at order $m^0$; an odd piece
contributes only at order $1/m$, through the $\beta\varepsilon_{\text{odd}}^2/2m$
term. This is why coefficients coupling through odd operators (such as
$b_0$ and $H_{0i}$) are systematically subleading relative to those
coupling through even operators (such as $b_i$ and $H_{ij}$). Projecting
the block-diagonal result onto the upper block gives the matter-sector
Hamiltonian directly; the antiparticle sector requires the separate
hole-theory reinterpretation of the lower (negative-energy) block, not a
direct charge-conjugation transformation of the field-theoretic coupling
itself, since the relevant bilinear $\bar\psi\gamma_5\gamma^\mu\psi$ is
even under charge conjugation for every $\mu$ and so cannot alone account
for the sign difference between the two sectors (Chapter~4).

\subsection{CPT Symmetry and the Asymmetry Parameter}
\label{sec:asymmetry-def}

CPT symmetry, invariance under the combined transformation of charge
conjugation, parity, and time reversal, is a theorem of local,
Lorentz-invariant quantum field theory (Colladay and Kostelecký, 1997). To
quantify possible CPT asymmetries in exotic interactions, this thesis
defines the dimensionless asymmetry parameter
\begin{equation}
  A_{\alpha}
  = \frac{g_{\alpha}^{f} - g_{\alpha}^{\bar{f}}}
         {g_{\alpha}^{f} + g_{\alpha}^{\bar{f}}},
  \qquad \alpha\in\{s,\,p,\,V,\,A\},
  \label{eq:asymmetry}
\end{equation}
where $g_{\alpha}^{f}$ and $g_{\alpha}^{\bar f}$ are the effective coupling
constants for a fermion species and its antifermion in interaction channel
$\alpha$. Under exact CPT symmetry, $A_\alpha=0$ for every channel. As
Chapter~4 derives in detail, $g_\alpha^f$ and $g_\alpha^{\bar f}$ as
actually reported in the literature are one-sided experimental upper
bounds rather than signed measured couplings, which limits how directly a
non-zero $A_\alpha$ can be read as evidence of CPT violation.

\section{Review of Existing Knowledge}
\label{sec:review}

\subsection{Experimental Platforms for Exotic Spin-Dependent Interactions}

The search for exotic spin-dependent interactions spans a diverse range of
experimental platforms, each optimised for a different length scale,
particle species, and sensitivity regime. Nitrogen-vacancy (NV) centres in
diamond are sensitive solid-state quantum sensors for weak spin-dependent
fields at sub-micrometre scales, owing to long coherence times and optical
spin readout (Rong \textit{et al.}, 2018a, 2018b), and are effective for
short-range Yukawa-type potentials from $0.1\,\mu\mathrm{m}$ to
$50\,\mu\mathrm{m}$. Torsion-pendulum experiments provide macroscopic
sensitivity to weak spin-dependent and spin-independent forces over
mesoscopic distances, typically $10\,\mu\mathrm{m}$ to $10\,\mathrm{cm}$,
using polarised test masses against unpolarised attractors (Heckel
\textit{et al.}, 2008). Atomic magnetometers operating in the
spin-exchange relaxation-free regime achieve sub-femtotesla sensitivity to
exotic spin couplings and are especially relevant for constraining $b_\mu$
and $d_{\mu\nu}$ (Vasilakis \textit{et al.}, 2009). Antimatter experiments
at CERN provide a direct test of CPT symmetry: BASE measures the antiproton
magnetic moment to parts-per-billion precision (Smorra \textit{et al.},
2017), and ALPHA performs precision spectroscopy on trapped antihydrogen
(Ahmadi \textit{et al.}, 2017). Positronium, a purely leptonic bound
system free of hadronic uncertainties, provides sensitivity to
electron--positron spin couplings at atomic scales, with recent
improvements from hyperfine and fine-structure spectroscopy (Fadeev
\textit{et al.}, 2022; Adkins, Cassidy and Pérez-Ríos, 2022).

\subsection{Theoretical Development History}

Moody and Wilczek (1984) gave the first systematic treatment of long-range
spin-dependent forces mediated by axion-like particles, establishing the
Yukawa-type potential structure and the $g_sg_p$ coupling dependence that
remains standard today. Colladay and Kostelecký (1997, 1998) constructed
the SME as the most general effective field theory allowing Lorentz and CPT
violation, and Kostelecký and Lane (1999) derived the non-relativistic
limit of SME fermion-sector terms, providing the first formal bridge
toward experimentally accessible observables. Dobrescu and Mocioiu (2006)
then provided the complete classification of non-relativistic potentials
described in Section~\ref{sec:dm}. More recently, Fadeev \textit{et al.}
(2019, 2022) unified exotic spin-dependent interactions under a common
effective field theory language and extended the analysis to observables in
simple atomic systems, while Cong \textit{et al.} (2025) compiled the most
complete modern survey of experimental constraints and exotic interaction
classifications through 2024, without, however, constructing a systematic
SME--Dobrescu--Mocioiu mapping or a matter--antimatter comparison
framework.

\section{Problems and Inadequacies in Existing Knowledge}
\label{sec:gap}

Despite these theoretical and experimental advances, the field remains
fragmented across multiple formalisms and experimental reporting
conventions, which limits cross-platform comparability and prevents a
coherent programme of CPT tests using exotic interactions. The SME
provides a fully relativistic description through well-defined operator
coefficients, but these coefficients are not directly expressed in terms
of experimentally measurable potentials. The Dobrescu--Mocioiu framework
provides an experimentally convenient basis, but is phenomenological and
does not make its connection to relativistic coefficients explicit. Prior
attempts at bridging the two via Foldy--Wouthuysen reduction are
incomplete, typically restricted to a subset of SME coefficients or a
limited class of potentials; a fully general mapping between the complete
SME fermion sector and the full $V_1$--$V_{16}$ basis has not previously
been established. Existing experimental compilations, including the recent
and extensive one of Cong \textit{et al.} (2025), are not expressed within
a unified parameter space that allows direct comparison across particle
species, interaction types, and experimental platforms. Most critically, no
existing framework defines a universal CPT asymmetry parameter across all
exotic spin-dependent interaction channels, nor has such a parameter been
systematically evaluated against compiled experimental data. This prevents
a quantitative, model-independent assessment of CPT violation in
spin-dependent forces, and is the gap this thesis addresses.

\section{Limitations of Existing Approaches}
\label{sec:existing-limitations}

Three specific limitations of the existing literature bear directly on the
work of this thesis. First, published compilations report constraints as
one-sided upper bounds, not signed central values with symmetric
uncertainties; any CPT-comparison framework built on such data inherits
this limitation and must, as shown in Chapter~4, treat it explicitly rather
than assume it away. Second, antimatter-sector measurements are far rarer
and far less standardised in units and reporting convention than
matter-sector measurements, which existing compilations do not
systematically quantify sector by sector or potential by potential. Third,
the SME--Dobrescu--Mocioiu correspondence has, prior to this work, only
been worked out piecemeal, coefficient by coefficient, in different papers
using different conventions for signs, index placement, and normalisation,
which makes an independent, internally consistent re-derivation, of the
kind undertaken in Chapters~3 and~4 of this thesis, necessary before a
reliable matching table can be trusted for quantitative use.
